\documentclass[a4paper,twoside]{article}

\usepackage{epsfig}
\usepackage{subfigure}
\usepackage{calc}
\usepackage{amssymb}
\usepackage{amstext}
\usepackage{amsmath}
\usepackage{amsthm}
\usepackage{multicol}
\usepackage{pslatex}
%\usepackage{hyperref} % use this for interactive urls and citations
\usepackage{ICAIIT} 
\usepackage[shortlabels]{enumitem}


\setlist[itemize]{label=\rule[0.5ex]{0.6ex}{0.6ex}}
%ist[itemize]{leftmargin=0.9cm}
%\setlist[enumerate]{leftmargin=1cm}

%\setlist[itemize]{align=parleft,left=1em}
%\setlist[enumerate]{align=parleft,left=1em}

\setlist[enumerate,1]{leftmargin=1cm}
\setlist[itemize,1]{leftmargin=0.9cm}



\subfigtopskip=0pt
\subfigcapskip=0pt
\subfigbottomskip=0pt


\begin{document}
\title{Authors' Instructions} 

\author{\authorname{First Author Name Surname\sup{1}, Second Author Name Surname\sup{1} and Third Author Name Surname\sup{2}}
\affiliation{\sup{1}Institute of Problem Solving, XYZ University, My Street, MyTown, MyCountry}
\affiliation{\sup{2}Department of Computing, Main University, MySecondTown, MyCountry}
\email{\{f\_author, s\_author\}@ips.xyz.edu, t\_author@dc.mu.edu}
}

\keywords{The paper must have at least one keyword. This paragraph should be justified with a line space of exactly 11-point, special indentation of 2-centimeters and with a spacing before of exactly 48-points. The text should be set in 9-point font size and without the use of bold or italic font style. For more than one keyword, please use a comma as a separator. Keywords must be title cased.}

\abstract{The abstract should summarize the contents of the paper and should contain at least 150 and at most 300 words. It should be set in 9-point font size, justified and should have a hanging indent of 2-centimeter. There should be a space before of 12-point and after of 30-point. The abstract must include research actuality, aims of the paper, research method, author's results, a practical and theoretical relevance of results.}

\onecolumn \maketitle \normalsize \vfill

\section{\uppercase{Introduction}}
\label{sec:introduction}

\noindent Your paper will be part of the conference proceedings; therefore, we ask that authors follow the guidelines explained in this example, in order to achieve the highest quality possible. Be advised that papers in a technically unsuitable form will be returned for retyping. After returned the manuscript must be appropriately modified. 


Margins, column widths, line spacing, and type styles are built-in. Some components, such as multi-levelled equations, graphics, and tables are not prescribed, although the various table text styles are provided. The formatter will need to create these components, incorporating the applicable criteria that follow.

\section{\uppercase{Manuscript Preparation}}

\noindent We strongly encourage authors to use this document for the preparation of the camera-ready. Please follow the instructions closely in order to make the volume look as uniform as possible.


Please remember that all the papers must be in English and without orthographic errors.


Do not add any text to the headers (do not set running heads) and footers, not even page numbers, because text will be added electronically.


For a best viewing experience the used font must be Times New Roman, on a Macintosh use the font named times, except on special occasions, such as program code (Section \ref{subsubsec:program_code}).

\subsection{Page Setup}

\noindent The paper size must be set to A4 (210x297 mm). The document
margins must be the following:

\begin{itemize}
    \item Top: 3,3 cm;
    \item Bottom: 4,2 cm;
    \item Left: 2,6 cm;
    \item Right: 2,6 cm.
\end{itemize}

It is advisable to keep all the given values.


Regarding the page layout, authors should set the Section Start to Continuous with the vertical alignment to the top and the following header and footer:

\begin{itemize}
    \item Header: 1,4 cm;
    \item Footer: 2,5 cm.
\end{itemize}

Any text or material outside the aforementioned margins will not be printed.

\subsection{First Section}

\noindent This section must be in one column.

\vfill
\subsubsection{Title and Subtitle}

\noindent Use the command \textit{$\backslash$title} and follow the given structure in "example.tex". The title and subtitle must be with initial letters
capitalized (titlecased). If no subtitle is required, please remove the corresponding \textit{$\backslash$subtitle} command. In the title or subtitle, words like "is", "or", "then", etc. should not be capitalized unless they are the first word of the subtitle. No formulas or special characters of any form or language are allowed in the title or subtitle.

\subsubsection{Authors and Affiliations}

\noindent Use the command \textit{$\backslash$author} and follow the given structure in "example.tex".

\subsubsection{Keywords}

\noindent Use the command \textit{$\backslash$keywords} and follow the given structure in "example.tex". Each paper must have at least one keyword. If more than one is specified, please use a comma as a separator. The sentence must end with a period.

\subsubsection{Abstract}

\noindent Use the command \textit{$\backslash$abstract} and follow the given structure in "example.tex".
Each paper must have an abstract up to 300 words. The sentence
must end with a period.

\subsection{Second Section}

\noindent Files "example.tex" and "example.bib" show how to create a paper with a 
corresponding list of references.

This section must be in two columns.

Each column must be 7,5-centimeter wide with a column spacing
of 0,8-centimeter.

The section text must be set to 10-point.

Section, subsection and sub-subsection first paragraph should not
have the first line indent.

To remove the paragraph indentation (only necessary for the
sections), use the command \textit{$\backslash$noindent} before the
paragraph first word.

If you use other style files (.sty) you MUST include them in the
final manuscript zip file.

\subsubsection{Section Titles}

\noindent The heading of a section title should be in all-capitals.

Example: \textit{$\backslash$section\{FIRST TITLE\}}

\vfill
\subsubsection{Subsection Titles}

\noindent The heading of a subsection title must be with initial letters
capitalized (titlecased).

Words like "is", "or", "then", etc. should not be capitalized unless
they are the first word of the subsection title.

Example: \textit{$\backslash$subsection\{First Subtitle\}}

\subsubsection{Sub-Subsection Titles}

\noindent The heading of a sub subsection title should be with initial letters
capitalized (titlecased).

Words like "is", "or", "then", etc should not be capitalized unless
they are the first word of the sub subsection title.

Example: \textit{$\backslash$subsubsection\{First Subsubtitle\}}

\subsubsection{Tables}

\noindent Tables must appear inside the designated margins or they may span
the two columns.

Tables in two columns must be positioned at the top or bottom of the
page within the given margins. To span a table in two columns please add an asterisk (*) to the table \textit{begin} and \textit{end} command.

Example: \textit{$\backslash$begin\{table*\}}

\hspace*{1.5cm}\textit{$\backslash$end\{table*\}}\\

Tables should be centered and should always have a caption
positioned above it. The font size to use is 9-point. No bold or
italic font style should be used.

The final sentence of a caption should end with a period.

\begin{table}[h]
\caption{This caption has one line so it is
centered.}\label{tab:example1} \centering
\begin{tabular}{|c|c|}
  \hline
  Example column 1 & Example column 2 \\
  \hline
  Example text 1 & Example text 2 \\
  \hline
\end{tabular}
\end{table}

\begin{table}[h]
\caption{This caption has more than one line so it has to be
justified.}\label{tab:example2} \centering
\begin{tabular}{|c|c|}
  \hline
  Example column 1 & Example column 2 \\
  \hline
  Example text 1 & Example text 2 \\
  \hline
\end{tabular}
\end{table}

Please note that the word "Table" is spelled out.


\subsubsection{Figures}

\noindent Please produce your figures electronically, and integrate them into
your document and zip file.

Check that in line drawings, lines are not interrupted and have a
constant width. Grids and details within the figures must be clearly
readable and may not be written one on top of the other.

Figure resolution should be at least 300 dpi.

Figures must appear inside the designated margins or they may span
the two columns.

Figures in two columns must be positioned at the top or bottom of
the page within the given margins. To span a figure in two columns please add an asterisk (*) to the figure \textit{begin} and \textit{end} command.

Example: \textit{$\backslash$begin\{figure*\}}

\hspace*{1.5cm}\textit{$\backslash$end\{figure*\}}

Figures should be centered and should always have a caption
positioned under it. The font size to use is 9-point. No bold or
italic font style should be used.

\begin{figure}[!h]
  %\vspace{-0.2cm}
  \centering
   \includegraphics[width=\linewidth]{ICAIIT.png}
  \caption{This caption has one line so it is centered.}
  \label{fig:example1}
 \end{figure}

\begin{figure}[!h]
  \vspace{-0.2cm}
  \centering
   \includegraphics[width=\linewidth]{ICAIIT.png}
  \caption{This caption has more than one line so it has to be set to justify.}
  \label{fig:example2}
  \vspace{-0.1cm}
\end{figure}

The final sentence of a caption should end with a period.



Please note that the word "Figure" is spelled out.

\subsubsection{Equations}

\noindent Equations should be placed on a separate line, numbered and
centered.\\The numbers accorded to equations should appear in
consecutive order inside each section or within the contribution,
with the number enclosed in brackets and justified to the right,
starting with the number 1.

Example:

\begin{equation}\label{eq1}
    a=b+c
\end{equation}

\subsubsection{List}\label{subsubsec:list}

\noindent The following no numbered  list style:

\begin{itemize}
    \item 1;
    \item 2;
    \item 3;
    \item 4.
\end{itemize}

or numbered list style:

\begin{enumerate}
    \item 1;
    \item 2;
    \item 3;
    \item 4.
\end{enumerate}

If use multilevel list:

\begin{enumerate}[A)]
    \item 1 level;
    \item 1 level;
        \begin{enumerate}[1)]
            \item 2 level;
            \item 2 level;
            \begin{itemize}
                \item 3 level;
                \item 3 level
            \end{itemize}
        \end{enumerate}
\end{enumerate}

\subsubsection{Program Code}\label{subsubsec:program_code}

\noindent Program listing or program commands in text should be set in
typewriter form such as Courier New.

Example of a Computer Program in Pascal:

\begin{small}
\begin{verbatim}
 Begin
     Writeln('Hello World!!');
 End.
\end{verbatim}
\end{small}


The text must be aligned to the left and in 9-point type.

\vfill
\subsubsection{Reference Text and Citations}

\noindent Papers must have between 10 to 30 references. Try to exclude links to online resources and national non-translated publications. At least the half of the references must be published within the last 5-7 years are included in the WoS and SCOPUS provided DOI number. The template will number citations consecutively within brackets \cite{eason1955certain} (IEEE style, see more - https://www.zotero.org/styles/ieee). The sentence punctuation follows the bracket \cite{maxwell1873treatise, salter2011geographical, kampmann2019using}. Refer simply to the reference number, as in \cite{jacobs1963fine} --- do not use “Ref. \cite{maxwell1873treatise}” or “reference \cite{maxwell1873treatise}” except at the beginning of a sentence: “Reference \cite{maxwell1873treatise} was the first ...”


Number footnotes separately in superscripts. Place the actual footnote at the bottom of the column in which it was cited. Do not put footnotes in the reference list. Use letters for table footnotes.


Unless there are six authors or more give all authors’ names; do not use “et al.”. Papers that have not been published, even if they have been submitted for publication, should be cited as “unpublished”. Papers that have been accepted for publication should be cited as “in press”. Capitalize only the first word in a paper title, except for proper nouns and element symbols.


For papers published in translation journals, please give the English citation \cite{yorozu1987electron}.


Unless there are six authors or more give all authors’ names; do not use “et al.”.

\section{\uppercase{Copyright Form}}

\noindent For the mutual benefit and protection of Authors and
Publishers, it is necessary that Authors provide formal written
Consent to Publish and Transfer of Copyright before publication of
the Book. The signed Consent ensures that the publisher has the
Author's authorization to publish the Contribution.

The copyright form is located on the authors' reserved area.

The form should be completed and signed by one author on
behalf of all the other authors.

\section{\uppercase{Conclusions}}
\label{sec:conclusion}

\noindent Please note that ONLY the files required to compile your paper should be submitted. Previous versions or examples MUST be removed from the compilation directory before submission.

We hope you find the information in this template useful in the preparation of your submission.

\section*{\uppercase{Acknowledgements}}

\noindent If any, should be placed before the references section
without numbering. To do so please use the following command:
\textit{$\backslash$section*\{ACKNOWLEDGEMENTS\}}




\bibliographystyle{IEEEtran}
{\small
\bibliography{example}}

%\begin{thebibliography}{999}
%\bibitem[1]G. Eason, B. Noble, and I.N. Sneddon, “On certain integrals of Lipschitz-Hankel type involving products of Bessel functions,” Phil. Trans. Roy.
%\bibitem[2]J. Clerk Maxwell, A Treatise on Electricity and Magnetism, 3rd ed., vol. 2. Oxford: Clarendon, 1892, pp.68-73.
%\end{thebibliography}


\section*{\uppercase{Appendix}}

\noindent If any, the appendix should appear directly after the
references without numbering, and not on a new page. To do so please use the following command:
\textit{$\backslash$section*\{APPENDIX\}}

\vfill
\end{document}

